\documentclass[12pt]{article}
\usepackage[utf8]{inputenc}
\usepackage[T1]{fontenc}
\usepackage{hyperref}
\title{Pecheneg Language Resources for K-12 Curriculum Development}
\author{Compiled by Claude Sonnet 5.5 (claude-sonnet-5-5) for Andr\'as Kornai}
\date{}
\begin{document}
\maketitle

\section{Introduction}
Pecheneg (Patzinak; ISO 639-3 xpc, Glottocode pech1242) was the language of the medieval Pecheneg confederation of the Eurasian steppe. It is extinct. It is generally classed as Turkic, probably Oghuz, but it is poorly attested. The evidence is mostly names and a few glosses recorded in Greek, chiefly by Constantine Porphyrogenitus, and the Turkic reading is contested. A 2020 analysis by Mudrak argues that many names fit Ossetic origins better. No native-language school curriculum is possible, and the searches found no teaching materials of any kind.

\section{Primary Educational Resources}
None found. There are no primers, textbooks or school programmes.

\section{Linguistic and Reference Materials}
\begin{itemize}
\item Pechenegs (Wikipedia), a historical and linguistic overview: \url{https://en.wikipedia.org/wiki/Pechenegs}
\item Gábor Vőrős (2002), ``Relics of the Pecheneg language in the works of Constantine,'' listed in Glottolog (I did not read the text): \url{https://glottolog.org/resource/reference/id/hh:ld:Voros:Pecheneg}
\item Mudrak's analysis of Pecheneg names (abstract): \url{https://istina.ips.ac.ru/publications/article/330365212/}
\item ``Pecheneg names and titles'' (in Russian), a record in a Ukrainian repository: \url{https://nasplib.isofts.kiev.ua/handle/123456789/172374}
\item Cambridge University Press book chapter, returned by the search and not examined: \url{https://www.cambridge.org/core/product/identifier/CBO9780511815638A014/type/BOOK_PART}
\end{itemize}

\section{Religious and Community Resources}
None found. There is no known Pecheneg-speaking community and no scripture translation.

\section{Digital and Audio Resources}
None found.

\section{Curriculum Development Context}
The language is known only from names and a few glosses, and its genetic affiliation is disputed. The most that is defensible for schools is a history and source-criticism unit. It could cover who the Pechenegs were, how the Greek and Arabic sources record their language, and why scholars disagree about whether it was Turkic. Pecheneg itself cannot be the language of instruction.
\end{document}